% ===============================
% MAT221 -- Real Analysis
% Measure Theory Crash Course
% ===============================

\documentclass[12pt,a4paper]{article}

% ===============================
% Metadata
% ===============================
\newcommand*{\CourseCode}{MAT221}
\newcommand*{\CourseTitle}{Real Analysis}
\newcommand*{\Chapter}{Measure Theory Crash Course}
\newcommand*{\Topics}{$\sigma$-Algebras, Measures, Outer Measure, Lebesgue Measure, and Non-Measurable Sets}
\newcommand*{\DocDate}{September 01, 2026}

% ===============================
% Header
% ===============================
\input{header}

% References tailored to this lecture (without changing the shared header file).
\renewcommand{\ReferenceList}{%
\begin{enumerate}[leftmargin=2em,label={[\arabic{*}]}]
    \item Sheldon Axler, \textit{Measure, Integration \& Real Analysis}, Springer, 2020.
    \item H. L. Royden and P. M. Fitzpatrick, \textit{Real Analysis}, 4th Edition, Pearson, 2010.
    \item Gerald B. Folland, \textit{Real Analysis}, 2nd Edition, Wiley, 1999.
\end{enumerate}
}

% ===============================
% Document
% ===============================
\begin{document}

\FrontPage
\Large

% ================================================================
% ================================================================
\begin{heading}
    Why Measure Theory?
\end{heading}

\clearpage

% ================================================================
\begin{heading}
    1. $\sigma$-Algebras: The Measurable Sets
\end{heading}

\vspace{10pt}

\begin{definition}[$\sigma$-Algebra]
    Let $X$ be a set. A \textbf{$\sigma$-algebra} $\mathcal{A}$ on $X$ is a collection of subsets of $X$ such that:
    \begin{enumerate}
        \item $X\in\mathcal{A}$;
        \item if $A\in\mathcal{A}$, then $A^c=X\setminus A\in\mathcal{A}$;
        \item if $A_1,A_2,\ldots\in\mathcal{A}$, then
        $\displaystyle \bigcup_{n=1}^{\infty}A_n\in\mathcal{A}$.
    \end{enumerate}
    The pair $(X,\mathcal{A})$ is called a \textbf{measurable space}, and the members of $\mathcal{A}$ are called \textbf{measurable sets}.
\end{definition}

\clearpage

\begin{definition}[Basic Examples]
    For every set $X$:
    \begin{enumerate}
        \item $\{\emptyset,X\}$ is the \textbf{trivial $\sigma$-algebra};
        \item $\mathcal{P}(X)$ is the \textbf{power-set $\sigma$-algebra};
        \item if $X={1,2,3,4}$ and $\mathcal{A}=\{\emptyset,\{1,2\},\{3,4\},X\}$, then $\mathcal{A}$ is a $\sigma$-algebra on $X$.
    \end{enumerate}
\end{definition}

\clearpage

\begin{theorem}[Consequences of the Axioms]
    If $\mathcal{A}$ is a $\sigma$-algebra on $X$, then:
    \begin{enumerate}
        \item $\emptyset\in\mathcal{A}$;
        \item $\mathcal{A}$ is closed under finite unions and finite intersections;
        \item if $A_1,A_2,\ldots\in\mathcal{A}$, then
        $\displaystyle \bigcap_{n=1}^{\infty}A_n\in\mathcal{A}$;
        \item if $A,B\in\mathcal{A}$, then $A\setminus B\in\mathcal{A}$.
    \end{enumerate}
\end{theorem}

\clearpage

\begin{theorem}[Intersection of $\sigma$-Algebras]
    Let $\{\mathcal{A}_i\}_{i\in I}$ be any nonempty collection of $\sigma$-algebras on the same set $X$. Then
    \[
        \bigcap_{i\in I}\mathcal{A}_i
    \]
    is a $\sigma$-algebra on $X$.
\end{theorem}

\clearpage

\begin{definition}[Generated $\sigma$-Algebra]
    If $\mathcal{C}\subseteq\mathcal{P}(X)$, the \textbf{$\sigma$-algebra generated by $\mathcal{C}$} is
    \[
        \sigma(\mathcal{C})
        =\bigcap\{\mathcal{A}:\mathcal{A}\text{ is a $\sigma$-algebra on }X
        \text{ and }\mathcal{C}\subseteq\mathcal{A}\}.
    \]
    Thus $\sigma(\mathcal{C})$ is the smallest $\sigma$-algebra containing $\mathcal{C}$.
\end{definition}

\clearpage

\begin{definition}[Borel $\sigma$-Algebra]
    The \textbf{Borel $\sigma$-algebra} on $\mathbb{R}$ is
    \[
        \mathcal{B}(\mathbb{R})
        =\sigma\bigl(\{(a,b):a<b,\ a,b\in\mathbb{R}\}\bigr).
    \]
    Its elements are called \textbf{Borel sets}.
\end{definition}

\clearpage

\begin{problem}[Cardinality of the Borel $\sigma$-Algebra]
    Show that $|\mathcal{B}(\mathbb{R})|=\mathfrak{c}$. Hint: use the countable family of rational-endpoint intervals for the upper bound, and note that all singletons are Borel for the lower bound.
\end{problem}

\clearpage

% ================================================================
\begin{heading}
    2. Measures and Measure Spaces
\end{heading}

\vspace{10pt}

\begin{definition}[Measure]
    Let $(X,\mathcal{A})$ be a measurable space. A \textbf{measure} is a function
    \[
        \mu:\mathcal{A}\longrightarrow[0,\infty]
    \]
    such that:
    \begin{enumerate}
        \item $\mu(\emptyset)=0$;
        \item if $A_1,A_2,\ldots\in\mathcal{A}$ are pairwise disjoint, then
        \[
            \mu\!\left(\bigcup_{n=1}^{\infty}A_n\right)
            =\sum_{n=1}^{\infty}\mu(A_n).
        \]
    \end{enumerate}
    The triple $(X,\mathcal{A},\mu)$ is called a \textbf{measure space}.
\end{definition}

\clearpage

\begin{problem}[Counting Measure]
    On $(X,\mathcal{P}(X))$, define
    \[
        \#(A)=
        \begin{cases}
            \text{the number of elements of $A$},&A\text{ is finite},\\
            \infty,&A\text{ is infinite}.
        \end{cases}
    \]
    Show that $\#$ is a measure. Determine when it is finite and when it is $\sigma$-finite.
\end{problem}

\clearpage

\begin{problem}[Dirac Measure]
    Let $(X,\mathcal{A})$ be a measurable space and fix $x_0\in X$. Define
    \[
        \delta_{x_0}(A)=
        \begin{cases}
            1,&x_0\in A,\\
            0,&x_0\notin A.
        \end{cases}
    \]
    Show that $\delta_{x_0}$ is a finite measure on $(X,\mathcal{A})$.
\end{problem}

\clearpage

\begin{problem}[Probability Measure]
    Let $(X,\mathcal{A})$ be a measurable space. A measure $\mu$ is a \textbf{probability measure} if $\mu(X)=1$. Where we assume that $\mu(X)=1$, we can interpret $\mu(A)$ as the probability of the event $A\in\mathcal{A}$.
\end{problem}

\clearpage

\begin{problem}[Borel Measure]
    Let $\mathcal{A}=\mathcal{B}(\mathbb{R})$. There is a unique measure $\lambda$ on $(\mathbb{R},\mathcal{B}(\mathbb{R}))$ such that
    \[
        \lambda((a,b))=b-a\qquad(a<b).
    \]
    This is called \textbf{Borel measure}. Also for half open or closed intervals,
    \[\lambda((a,b])=\lambda([a,b))=\lambda([a,b])=b-a.\]
\end{problem}

\clearpage

\begin{theorem}[Basic Properties of Measures]
    Let $(X,\mathcal{A},\mu)$ be a measure space. Then:
    \begin{enumerate}
        \item \textbf{Monotonicity:} if $A\subseteq B$, then $\mu(A)\leq\mu(B)$.
        \item \textbf{Difference formula:} if $A\subseteq B$ and $\mu(A)<\infty$, then
        \[
            \mu(B\setminus A)=\mu(B)-\mu(A).
        \]
        \item \textbf{Countable subadditivity:}
        \[
            \mu\!\left(\bigcup_{n=1}^{\infty}A_n\right)
            \leq\sum_{n=1}^{\infty}\mu(A_n).
        \]
        \item \textbf{Continuity from below:} if $A_n\uparrow A$, then
        \[
            \mu(A)=\lim_{n\to\infty}\mu(A_n).
        \]
        \item \textbf{Continuity from above:} if $A_n\downarrow A$ and $\mu(A_1)<\infty$, then
        \[
            \mu(A)=\lim_{n\to\infty}\mu(A_n).
        \]
    \end{enumerate}
\end{theorem}

\clearpage

% ================================================================
\begin{heading}
    3. Outer Measure: Measuring Every Subset
\end{heading}

\vspace{10pt}

\begin{definition}[Outer Measure]
    An \textbf{outer measure} on a set $X$ is a function
    \[
        \mu^*:\mathcal{P}(X)\longrightarrow[0,\infty]
    \]
    such that:
    \begin{enumerate}
        \item $\mu^*(\emptyset)=0$;
        \item if $A\subseteq B\subseteq X$, then $\mu^*(A)\leq\mu^*(B)$;
        \item for every sequence $A_1,A_2,\ldots\subseteq X$,
        \[
            \mu^*\!\left(\bigcup_{n=1}^{\infty}A_n\right)
            \leq\sum_{n=1}^{\infty}\mu^*(A_n).
        \]
    \end{enumerate}
    Unlike a measure, an outer measure is defined on \emph{all} subsets of $X$, but it need not be additive on all of them.
\end{definition}

\clearpage

\begin{definition}[Length of an Interval]
    For any bounded interval $I$ with endpoints $a\leq b$, define
    \[
        \ell(I)=b-a,
    \]
    regardless of whether either endpoint belongs to $I$. Set $\ell(\emptyset)=0$.
\end{definition}

\vspace{30pt}

\begin{definition}[Lebesgue Outer Measure on $\mathbb{R}$]
    For every $A\subseteq\mathbb{R}$, define
    \[
        \lambda^*(A)=
        \inf\left\{
        \sum_{n=1}^{\infty}\ell(I_n):
        A\subseteq\bigcup_{n=1}^{\infty}I_n,
        \ I_n\text{ bounded open intervals}
        \right\}.
    \]
    We use the convention $\inf\emptyset=\infty$.
\end{definition}


\clearpage

\begin{theorem}[Properties of Lebesgue Outer Measure]
    For all $A\subseteq\mathbb{R}$ and $t\in\mathbb{R}$:
    \begin{enumerate}
        \item if $I$ is a bounded interval, then $\lambda^*(I)=\ell(I)$;
        \item $\lambda^*(A+t)=\lambda^*(A)$, where
        $A+t=\{x+t:x\in A\}$;
        \item if $A$ is countable, then $\lambda^*(A)=0$.
    \end{enumerate}
\end{theorem}

\clearpage

\begin{problem}[Countable Sets Have Outer Measure zero]
    If $A=\{a_1,a_2,\ldots\}$ and $\varepsilon>0$, cover $a_n$ by an interval of length $\varepsilon/2^n$. Use this to prove $\lambda^*(A)=0$.
\end{problem}

\clearpage

\begin{problem}[Outer Measure Need Not Be Additive]
    Explain why we cannot yet call $\lambda^*$ a measure on $\mathcal{P}(\mathbb{R})$. Which axiom is missing?
\end{problem}

\clearpage

% ================================================================
\begin{heading}
    4. Carath\'eodory Measurability
\end{heading}

\vspace{10pt}

\begin{definition}[Carath\'eodory Criterion]
    Let $\mu^*$ be an outer measure on $X$. A set $E\subseteq X$ is \textbf{$\mu^*$-measurable} if, for every $A\subseteq X$,
    \[
        \mu^*(A)=\mu^*(A\cap E)+\mu^*(A\setminus E).
    \]
    Thus a measurable set splits every test set $A$ without creating or losing outer measure.
\end{definition}

\clearpage

\begin{theorem}[Carath\'eodory's Theorem]
    Let $\mu^*$ be an outer measure on $X$, and let
    \[
        \mathcal{M}_{\mu^*}
        =\{E\subseteq X:E\text{ is $\mu^*$-measurable}\}.
    \]
    Then $\mathcal{M}_{\mu^*}$ is a $\sigma$-algebra, and the restriction
    $\mu^*|_{\mathcal{M}_{\mu^*}}$ is a complete measure.
\end{theorem}

\clearpage

% ================================================================
\begin{heading}
    5. Lebesgue Measure
\end{heading}

\vspace{10pt}

\begin{definition}[Lebesgue Measurable Sets and Lebesgue Measure]
    Apply Carath\'eodory's construction to Lebesgue outer measure $\lambda^*$. The resulting $\sigma$-algebra is denoted by $\mathcal{L}(\mathbb{R})$ and its members are the \textbf{Lebesgue measurable sets}. The \textbf{Lebesgue measure} is
    \[
        \lambda=\lambda^*|_{\mathcal{L}(\mathbb{R})}.
    \]
    Hence, for $E\in\mathcal{L}(\mathbb{R})$,
    \[
        \lambda(E)=\lambda^*(E).
    \]
\end{definition}

\clearpage

\begin{theorem}[Borel Sets Versus Lebesgue Measurable Sets]
    Every Borel set is Lebesgue measurable, and
    \[
        \mathcal{B}(\mathbb{R})\subsetneq\mathcal{L}(\mathbb{R})
        \subsetneq\mathcal{P}(\mathbb{R}).
    \]
    The restriction $\lambda|_{\mathcal{B}(\mathbb{R})}$ is called \textbf{Borel Lebesgue measure}. It is the unique measure on $\mathcal{B}(\mathbb{R})$ satisfying
    \[
        \lambda((a,b])=b-a\qquad(a<b).
    \]
\end{theorem}

\clearpage

\begin{definition}[Complete Measure Space]
    A measure space $(X,\mathcal{A},\mu)$ is \textbf{complete} if whenever $N\in\mathcal{A}$, $\mu(N)=0$, and $S\subseteq N$, then $S\in\mathcal{A}$ (and necessarily $\mu(S)=0$).
\end{definition}

\vspace{30pt}

\begin{theorem}[Completion of Borel Lebesgue Measure]
    The Lebesgue $\sigma$-algebra is the completion of the Borel $\sigma$-algebra with respect to Borel Lebesgue measure. Equivalently, $E\subseteq\mathbb{R}$ is Lebesgue measurable if and only if there exist a Borel set $B$ and a subset $N$ of a Borel null set such that
    \[
        E=B\cup N.
    \]
\end{theorem}

\clearpage

\begin{theorem}[Principal Properties of Lebesgue Measure]
    Lebesgue measure is complete, translation invariant, and countably additive. In particular, if $E\in\mathcal{L}(\mathbb{R})$ and $t\in\mathbb{R}$, then $E+t$ is measurable and
    \[
        \lambda(E+t)=\lambda(E).
    \]
\end{theorem}

\clearpage

\begin{definition}[Null Set]
    A Lebesgue measurable set $N\subseteq\mathbb{R}$ is a \textbf{null set} if $\lambda(N)=0$.
\end{definition}

\vspace{50pt}

\begin{theorem}[Countable Sets Are Null]
    Every countable subset of $\mathbb{R}$ is Lebesgue measurable and has measure zero. In particular,
    \[
        \lambda(\{x\})=0
        \qquad\text{and}\qquad
        \lambda(\mathbb{Q})=0.
    \]
\end{theorem}


\clearpage

% ================================================================
\begin{heading}
    6. An Uncountable Null Set: The Cantor Set
\end{heading}

\vspace{10pt}

\begin{definition}[Cantor Set]
    Start with $C_0=[0,1]$. From every interval in $C_{n-1}$ remove its open middle third, obtaining $C_n$. The \textbf{Cantor set} is
    \[
        C=\bigcap_{n=0}^{\infty}C_n.
    \]
    At stage $n$, the set $C_n$ is a union of $2^n$ closed intervals, each of length $3^{-n}$.
\end{definition}

\clearpage

\begin{theorem}[Basic Properties of the Cantor Set]
    The Cantor set $C$ is closed, uncountable, has empty interior, and satisfies
    \[
        \lambda(C)=0.
    \]
\end{theorem}

\vspace{10pt}

\begin{problem}[Why Is the Cantor Set Null (Measure of zero)?]
    Since $C\subseteq C_n$, use
    \[
        \lambda(C_n)=2^n3^{-n}=\left(\frac{2}{3}\right)^n
    \]
    and let $n\to\infty$.
\end{problem}

\clearpage

\begin{problem}[Why Is the Cantor Set Uncountable?]
    Show that $x\in C$ precisely when $x$ has a ternary expansion using only the digits $0$ and $2$ (with the usual care at endpoints). Relate such expansions to infinite binary sequences.
\end{problem}

\clearpage

\begin{problem}[Cardinality of the Lebesgue $\sigma$-Algebra]
    Prove that
    \[
        |\mathcal{L}(\mathbb{R})|=2^{\mathfrak{c}}.
    \]
    Hint: $|C|=\mathfrak{c}$ and every subset of the null set $C$ is Lebesgue measurable.
\end{problem}

\clearpage

\begin{heading}
    Comparison of Borel and Lebesgue Measurability
\end{heading}

\begin{theorem}[Strictness of the Borel--Lebesgue Inclusion]
    Since $|\mathcal{B}(\mathbb{R})|=\mathfrak{c}$ but
    $|\mathcal{L}(\mathbb{R})|=2^{\mathfrak{c}}$, there exist Lebesgue measurable sets that are not Borel sets.
\end{theorem}

\clearpage

% ================================================================
\begin{heading}
    7. Almost Everywhere
\end{heading}

\vspace{10pt}

\begin{definition}[Almost Everywhere]
    A property $P(x)$ holds \textbf{almost everywhere} (a.e.) on a measure space $(X,\mathcal{A},\mu)$ if there exists a measurable null set $N$ such that $P(x)$ holds for every $x\in X\setminus N$.

    If the failure set $\{x\in X:P(x)\text{ fails}\}$ is measurable, this is equivalent to saying that the failure set has measure zero.
\end{definition}

\clearpage

\begin{problem}[Equality Almost Everywhere]
    Let $f,g:X\to\mathbb{R}$. Write precisely what $f=g$ a.e. means. Show that ``equal almost everywhere'' is an equivalence relation (assuming the relevant sets are measurable).
\end{problem}

\clearpage

% ================================================================
\begin{heading}
    8. Not Every Subset Is Lebesgue Measurable
\end{heading}

\vspace{10pt}

\begin{definition}[Vitali Equivalence Relation]
    On $[0,1]$, define
    \[
        x\sim y\quad\Longleftrightarrow\quad x-y\in\mathbb{Q}.
    \]
    Using the axiom of choice, select one representative from each equivalence class. Any resulting set $V\subseteq[0,1]$ is called a \textbf{Vitali set}.
\end{definition}

\begin{theorem}[A Vitali Set Is Not Lebesgue Measurable]
    A Vitali set $V$ is not Lebesgue measurable.
\end{theorem}

\clearpage

\begin{problem}[Vitali Argument]
    Let $\mathbb{Q}\cap[-1,1]=\{q_1,q_2,\ldots\}$. Show that the sets $V+q_n$ are pairwise disjoint, and that
    \[
        [0,1]\subseteq\bigcup_{n=1}^{\infty}(V+q_n)\subseteq[-1,2].
    \]
    Assuming $V$ were measurable, use translation invariance and countable additivity to obtain a contradiction, considering separately the cases $\lambda(V)=0$ and $\lambda(V)>0$.
\end{problem}

\clearpage

% ================================================================
\begin{heading}
    9. Why Countable Additivity?
\end{heading}

\vspace{10pt}

\begin{problem}[Finite Additivity Is Too Weak]
    Countable additivity implies finite additivity, but not conversely. Explain why finite additivity alone does not force
    \[
        \mu\!\left(\bigcup_{n=1}^{\infty}A_n\right)
        =\sum_{n=1}^{\infty}\mu(A_n)
    \]
    for disjoint sets $A_n$, and why this would break continuity from below.
\end{problem}

\vspace{50pt}

\begin{problem}[Why Not Uncountable Additivity?]
    Every interval is the uncountable disjoint union of its singleton points. If a length measure were additive over arbitrary disjoint families, then
    \[
        \lambda([0,1])=\sum_{x\in[0,1]}\lambda(\{x\})=0,
    \]
    contradicting $\lambda([0,1])=1$. Explain why countable additivity avoids this contradiction.
\end{problem}

\clearpage

% ================================================================
\begin{heading}
    10. Higher-Dimensional Lebesgue Measure
\end{heading}

\vspace{10pt}

\begin{definition}[Volume of a Rectangle in $\mathbb{R}^n$]
    For a half-open rectangle
    \[
        R=\prod_{j=1}^{n}(a_j,b_j],
        \qquad
        \operatorname{vol}(R)=\prod_{j=1}^{n}(b_j-a_j).
    \]
    
\end{definition}

\clearpage

\begin{definition}[Lebesgue Outer Measure on $\mathbb{R}^n$]
    For a half-open rectangle
    \[
        R=\prod_{j=1}^{n}(a_j,b_j],
        \qquad
        \operatorname{vol}(R)=\prod_{j=1}^{n}(b_j-a_j).
    \]
    For $A\subseteq\mathbb{R}^n$, define
    \[
        \lambda_n^*(A)=
        \inf\left\{
        \sum_{k=1}^{\infty}\operatorname{vol}(R_k):
        A\subseteq\bigcup_{k=1}^{\infty}R_k
        \right\}.
    \]
    Carath\'eodory's criterion then produces the Lebesgue measurable subsets of $\mathbb{R}^n$ and Lebesgue measure $\lambda_n$.
\end{definition}

\clearpage

\begin{theorem}[Translation Invariance]
    If $E\subseteq\mathbb{R}^n$ is Lebesgue measurable and $t\in\mathbb{R}^n$, then $E+t$ is measurable and
    \[
        \lambda_n(E+t)=\lambda_n(E).
    \]
\end{theorem}

\clearpage

\begin{theorem}[What the Non-Measurability Results Actually Say]
    \begin{itemize}

    \item It is not possible to define the Lebesgue measure on \emph{all}
    subsets of $\mathbb{R}^n$ while preserving the usual geometric
    properties of measure.

    \item \textbf{Hausdorff's result (1914):}
    For every dimension $n \geq 1$, there is no countably additive measure
    defined on all subsets of $\mathbb{R}^n$ that
    \begin{itemize}
        \item is invariant under isometries, such as translations and rotations, and
        \item assigns measure $1$ to the unit cube.
    \end{itemize}

    \item Moreover, Hausdorff showed that for $n \geq 3$, no such
    \emph{finitely additive} measure exists.

    \item \textbf{Banach--Tarski paradox (1924):}
    If $n \geq 3$, a ball in $\mathbb{R}^n$ can be decomposed into finitely
    many pieces and, using only isometries, these pieces can be reassembled
    into a ball of any prescribed volume.

    \item Informally, this means that one could, in principle, decompose a
    small ball and reassemble the pieces into a much larger one---for example,
    ``reassemble a pea into the sun.''

    \item The construction of such pieces depends essentially on the
    \emph{axiom of choice}.

    \item \textbf{Banach's result (1923):}
    For $n=1$ and $n=2$, there do exist finitely additive,
    isometry-invariant extensions of Lebesgue measure to all subsets of
    $\mathbb{R}^n$.

    \item However, these extensions are \emph{not countably additive}.

    \item In dimensions $n \geq 3$, the obstruction is even stronger:
    finite additivity cannot be preserved under these requirements.

\end{itemize}
\end{theorem}


\clearpage

% ================================================================
\begin{heading}
    Summary: The Logical Flow
\end{heading}

\vspace{10pt}

\begin{enumerate}
    \item A $\sigma$-algebra specifies the sets on which countable set operations are safe.
    \item A measure is countably additive on a $\sigma$-algebra.
    \item Lebesgue outer measure assigns a preliminary size to every subset of $\mathbb{R}$.
    \item Carath\'eodory's criterion selects the sets on which that outer measure behaves additively.
    \item Restricting $\lambda^*$ to those sets gives complete Lebesgue measure.
    \item Every Borel set is Lebesgue measurable, but some Lebesgue measurable sets are not Borel.
    \item Null sets may be uncountable, while some subsets of $\mathbb{R}$ are not Lebesgue measurable at all.
\end{enumerate}

\begin{problem}[Concept Check]
    Explain the differences among the following objects:
    \[
        \mathcal{B}(\mathbb{R}),\qquad
        \mathcal{L}(\mathbb{R}),\qquad
        \mathcal{P}(\mathbb{R}),\qquad
        \lambda^*,\qquad
        \lambda.
    \]
    State the domain of each set function and the inclusion relations among the three collections of sets.
\end{problem}

\clearpage

\EndPage

\end{document}
